(1) 单振子能级 $E_n=\hbar\omega_E(n+1/2)$。记 $x=\hbar\omega_E/(k_BT)$，正则配分函数由几何级数得到 $z=e^{-x/2}/(1-e^{-x})$，$3N$ 个独立振子的总配分函数为 $Z=z^{3N}$，这里未考虑其相互作用。
(2) 热内能来自 $\ln Z$ 对逆温的导数，除去恒定零点能有 $U_{\rm th}=3N\hbar\omega_E/(e^x-1)$。对温度求导得到 $C_V=3Nk_Bx^2e^x/(e^x-1)^2$，在正温度区间显然非负。
(3) 低温 $x=\Theta_E/T\gg1$ 时，分母 $(e^x-1)^2\simeq e^{2x}$，因此 $C_V\simeq3Nk_B(\Theta_E/T)^2e^{-\Theta_E/T}$。这一指数形式来自全部振子具有相同非零激发能隙。
(4) 高温 $x\ll1$，使用 $e^x-1\simeq x$ 得 $C_V\to3Nk_B$。Debye固体低温含频率可趋零的声学模，热容按 $T^3$ 衰减，而Einstein模型低温是指数衰减；两模型在高温都回到每原子 $3k_B$。
